\chapter{Interfaces as Cognitive Prosthetics}

\chapterepigraph{%
  A good interface does not translate your intentions into machine actions.\\
  It extends the space of intentions you can have.
}{Flyxion, \textit{Semantic Infrastructure}}

\sidenote{Connects to\\projection-based\\consciousness (Ch.\,22)\\and constraint\\spaces}

The previous three chapters developed the theory of knowledge navigation.
This chapter examines the interfaces through which humans navigate:
not as passive tools that execute commands but as cognitive prosthetics
that extend the agent's projection structure.

\section{Tools That Extend Memory}

The constraint-preserving memory chain of Chapter~23 is biological:
it lives in neural tissue and dies with the agent. External tools extend
this chain into physical space, where it can survive the agent and be
accessed by others.

\begin{definition}{Memory Prosthetic}{memory-prosthetic}
  A \emph{memory prosthetic} is any external artifact that encodes
  constraint modifications $\Delta\mathcal{C}$ in a retrievable form.
  Its effectiveness is measured by:
  \begin{itemize}
    \item \textbf{Encoding fidelity:} how accurately it represents
          $\Delta\mathcal{C}$.
    \item \textbf{Retrieval efficiency:} how quickly the encoded
          constraint modification can be recovered and applied.
    \item \textbf{Durability:} how long the encoding persists.
    \item \textbf{Portability:} how easily the encoding can be
          transferred to new contexts.
  \end{itemize}
\end{definition}

Writing is the foundational memory prosthetic: high encoding fidelity,
moderate retrieval efficiency (requires reading), high durability,
high portability. Digital storage improves retrieval efficiency
dramatically but at the cost of technological dependence: the encoding
requires an interpreter (software, hardware) that may not persist.

\section{Tools That Extend Attention}

Attention is the agent's mechanism for allocating cognitive resources —
for deciding which features of the current observation to process in
depth. Interfaces that extend attention allow the agent to maintain
focus on more features simultaneously, or to maintain focus for longer.

\begin{definition}{Attention Prosthetic}{attention-prosthetic}
  An \emph{attention prosthetic} is an interface that reduces the
  cognitive cost of directing attention to specific features of the
  environment. Its effectiveness is measured by:
  \[
    E_\text{att} = \frac{\Delta I_\text{attended}}{
    \Delta \kappa_\text{cognitive}}
  \]
  — the increase in attended information per unit increase in cognitive
  constraint violation (cognitive load).
\end{definition}

Good visualization tools are attention prosthetics: they make structure
visible that would otherwise require sustained mental effort to construct.
A good graph reveals in a glance what a table of numbers reveals only
after minutes of computation.

\section{Tools That Extend Reasoning}

The most powerful cognitive prosthetics extend not just memory or
attention but reasoning itself: they change what inferences are
accessible, what concepts can be combined, what questions can be
asked.

Mathematical notation is the canonical example. A student who knows
Arabic numerals and positional notation can perform calculations that
a student who knows only Roman numerals cannot — not because they are
smarter, but because the notation makes certain operations cognitively
accessible that are inaccessible in the other system.

\begin{namedthm}{The Cognitive Prosthetic Principle}
  The cognitive prosthetics available to an agent determine the
  accessible region of their cognitive landscape. The accessibility
  entropy of an agent's cognitive state is:
  \[
    S_\text{cog}(A) = S_\text{bio}(A) + \sum_i S_\text{prosthetic}(P_i)
  \]
  — the sum of the agent's biological accessibility and the
  accessibility contributed by each prosthetic $P_i$ in their toolkit.
  Cognitive development is the acquisition of high-$S_\text{prosthetic}$
  tools.
\end{namedthm}

\begin{exercise}
  Evaluate the following tools as cognitive prosthetics: (a) a physical
  notebook, (b) a spreadsheet, (c) a programming language,
  (d) a diagram, (e) a conversational AI assistant. For each, estimate
  $E_\text{att}$ and $S_\text{prosthetic}$ informally. Which extends
  reasoning most? Which has the highest failure modes?
\end{exercise}

\begin{exercise}
  The abacus and the calculator both extend arithmetic reasoning.
  Compare them as cognitive prosthetics. Which has higher encoding
  fidelity? Which has higher retrieval efficiency? Which produces
  better long-term mathematical understanding? Explain the tradeoffs
  using the framework of this chapter.
\end{exercise}
